\section{模型评价与推广}
\subsection{模型特点}
本文的分层结构使决策边界可解释：问题一处理单区能力和离散组批，问题二在批次上叠加时限及设备/电池周转，问题三继承运输时间轴配置通信模式，问题四将冻结任务汇总到独立分区核算资源。峰值区间计数与库存缺口有明确的量纲和管理含义，且跨区任务通过连通组件约束保持完整。

\subsection{局限性}
本稿以已有参考解作为报告数据，没有重新运行求解器或独立复核参考解，因此所有数值应按初稿口径使用。问题三的通信模式及中继任务直接沿用参考结果，未在本轮复查各链路输入与实体资源占用。问题四固定运输/通信任务时序，不能解释联合重排后是否会有更小缺口；其峰值并发下界也不等同于考虑位置转场、组间协同后的实体设备规划。绘图所用汇总指标不包含随机天气和设备故障的不确定性。

\subsection{推广方向}
该框架可用于道路受阻地区的应急物资投送、灾时临时通信覆盖和多责任区资源配置。实际部署时，建议以实时气象、道路和通信监测更新输入，并采用滚动时域方式重排剩余任务；同时将中继机位置转场、频谱容量、充电功率与鲁棒备用量纳入统一资源模型。

\subsection{初稿数据来源声明}
问题一至问题三的图表与结果取自工作区 \texttt{results/reference/} 下的归档参考 JSON；问题四数值取自用户提供的更新版 \texttt{sections/8\_problem4.tex}。本文未重新求解 Q3/Q4，亦未审计或独立验证题解。文献检索因本机 SSL 证书校验失败未能完成，故本稿不列未经核实的外部文献。AI 辅助：使用 MathModel 内置 AI 助手（本会话模型标识 GPT-6 Astra）协助整理文本与绘图程序；MathModel/模型发布日期未由当前环境提供，提交前需按实际会话信息补充。
